\chapter{4位加法计数器——模$2^N$二进制计数器}

\section{应用统一框架}

\begin{enumerate}
  \item \textbf{模值N}：$N = 2^4 = 16$（4位二进制能表示0到15）
  \item \textbf{寄存器位宽k}：$k = 4$（Q3 Q2 Q1 Q0）
  \item \textbf{状态转移}：count $\leftarrow$ count + 1
  \item \textbf{回零条件}：当 count = 1111（15）时，下一状态自动归0
    \begin{itemize}
      \item 注意：4位二进制1111 + 1 = 10000（5位），但寄存器只有4位，最高位（进位）自然丢弃
      \item 所以\textbf{不需要显式清零}——自然溢出就是模16
    \end{itemize}
\end{enumerate}

\section{逐拍状态分析}

这是全书第一个需要逐拍分析的电路。请仔细看每个时钟周期内发生了什么。

\subsection{初始状态}

\clktick{0}（初始）：
\begin{itemize}
  \item Q3Q2Q1Q0 = 0000（状态0）
  \item 加法器计算：0000 + 1 = 0001
  \item D3D2D1D0 = 0001（准备好下一状态）
\end{itemize}

\subsection{时钟第1拍到第16拍}

\begin{beatanalysis}[4位计数器完整16拍]
\begin{center}
\begin{tabular}{|c|c|c|c|c|c|}
\hline
\textbf{时钟拍} & \textbf{上升沿前Q值} & \textbf{加法器输出}(Q+1) & \textbf{上升沿后Q值} & \textbf{十进制} & \textbf{进位} \\
\hline
1 & 0000 & 0001 & \textbf{0001} & 1 & 0 \\
2 & 0001 & 0010 & \textbf{0010} & 2 & 0 \\
3 & 0010 & 0011 & \textbf{0011} & 3 & 0 \\
4 & 0011 & 0100 & \textbf{0100} & 4 & 0 \\
5 & 0100 & 0101 & \textbf{0101} & 5 & 0 \\
6 & 0101 & 0110 & \textbf{0110} & 6 & 0 \\
7 & 0110 & 0111 & \textbf{0111} & 7 & 0 \\
8 & 0111 & 1000 & \textbf{1000} & 8 & 0 \\
9 & 1000 & 1001 & \textbf{1001} & 9 & 0 \\
10 & 1001 & 1010 & \textbf{1010} & 10 & 0 \\
11 & 1010 & 1011 & \textbf{1011} & 11 & 0 \\
12 & 1011 & 1100 & \textbf{1100} & 12 & 0 \\
13 & 1100 & 1101 & \textbf{1101} & 13 & 0 \\
14 & 1101 & 1110 & \textbf{1110} & 14 & 0 \\
15 & 1110 & 1111 & \textbf{1111} & 15 & 0 \\
16 & 1111 & 0000 & \textbf{0000} & 0 & \textbf{1} \\
\hline
\end{tabular}
\end{center}
\end{beatanalysis}

\textbf{关键观察}：
\begin{itemize}
  \item 第16拍：1111 + 1 = 10000，但在4位寄存器中只保存低4位（0000），进位自然的被丢弃
  \item 计数器\textbf{不需要任何清零逻辑}就能自动归零——因为4位自然溢出
  \item 这就是模16的本质：利用$2^N$的自然溢出
\end{itemize}

\section{为什么Q0翻转最快}

观察Q0（最低位）的变化：
$$\text{0, 1, 0, 1, 0, 1, 0, 1, ...}$$

Q0在每个时钟周期都翻转！频率 = $f_{clk} / 2$。

观察Q1的变化：
$$\text{0, 0, 1, 1, 0, 0, 1, 1, ...}$$

Q1每两个时钟周期翻转一次！频率 = $f_{clk} / 4$。

观察Q2的变化：频率 = $f_{clk} / 8$。

\textbf{规律}：$Q_k$ 的频率 = $f_{clk} / 2^{k+1}$。

\textbf{这就是分频器的原理——我们将在第四部分详细展开。}

\section{VHDL实现}

\begin{lstlisting}[caption={4位二进制计数器（模16）}]
library ieee;
use ieee.std_logic_1164.all;
use ieee.std_logic_unsigned.all;

entity counter_4bit is
  port (
    clk   : in  std_logic;
    rst_n : in  std_logic;
    count : out std_logic_vector(3 downto 0);
    carry : out std_logic  -- 计数一轮完成标志
  );
end counter_4bit;

architecture behavioral of counter_4bit is
  signal cnt : std_logic_vector(3 downto 0);
begin
  process(clk, rst_n)
  begin
    if rst_n = '0' then
      cnt <= (others => '0');
    elsif rising_edge(clk) then
      cnt <= cnt + 1;  -- 自然溢出，1111+1=0000
    end if;
  end process;

  count <= cnt;
  carry <= '1' when cnt = "1111" else '0';
end behavioral;
\end{lstlisting}

\section{内部寄存器变化过程}

将每个时钟上升沿后寄存器内容画出来：

\begin{center}
\begin{tikzpicture}[scale=1.0]
  \def\dx{0.7}      % 每拍宽度(cm)
  \def\hi{0.5}      % 高电平偏移
  \def\lo{0.05}     % 低电平偏移

  % === 纵坐标定义 ===
  \def\Yclk{1.4}
  \def\Yzero{2.8}
  \def\Yone{4.2}
  \def\Ytwo{5.6}
  \def\Ythree{7.0}

  % === 横轴刻度线 0-15 ===
  \foreach \n in {0,...,15} {
    \draw[gray!50, thin] ({\n*\dx}, 0.8) -- ({\n*\dx}, \Ythree+\hi+0.2);
    \node[font=\footnotesize] at ({\n*\dx}, 0.5) {\n};
  }
  % 拍16标记(回到0)
  \draw[gray!40, dashed, thin] ({16*\dx}, 0.8) -- ({16*\dx}, \Ythree+\hi+0.2);
  \node[font=\footnotesize, gray] at ({16*\dx}, 0.5) {0};

  % 横轴
  \draw[->] (-0.2, 0.8) -- ({16*\dx+0.3}, 0.8) node[right, font=\small] {时钟拍号};

  % === 纵向虚线(每个上升沿) ===
  \foreach \n in {0,...,15} {
    \draw[gray!25, dotted, thin] ({\n*\dx}, \Yclk-0.1) -- ({\n*\dx}, \Ythree+\hi+0.1);
  }

  % ====== CLK ======
  \draw[black, thick]
    (0, \Yclk+\lo) -- (0, \Yclk+\hi)
    \foreach \n in {0,...,15} {
      -- ({\n*\dx+0.5*\dx}, \Yclk+\hi) -- ({\n*\dx+0.5*\dx}, \Yclk+\lo) -- ({\n*\dx+\dx}, \Yclk+\lo)
    }
    -- ({16*\dx}, \Yclk+\hi);
  \node[font=\small] at ({16*\dx+0.5}, \Yclk+0.25) {clk};

  % ====== Q0 (blue): 0,1,0,1,... 每拍翻转 ======
  \draw[blue, very thick]
    (0,\Yzero+\lo) -- (0,\Yzero+\lo) -- (0,\Yzero+\hi) -- (1*\dx,\Yzero+\hi)
    -- (1*\dx,\Yzero+\lo) -- (2*\dx,\Yzero+\lo)
    -- (2*\dx,\Yzero+\hi) -- (3*\dx,\Yzero+\hi)
    -- (3*\dx,\Yzero+\lo) -- (4*\dx,\Yzero+\lo)
    -- (4*\dx,\Yzero+\hi) -- (5*\dx,\Yzero+\hi)
    -- (5*\dx,\Yzero+\lo) -- (6*\dx,\Yzero+\lo)
    -- (6*\dx,\Yzero+\hi) -- (7*\dx,\Yzero+\hi)
    -- (7*\dx,\Yzero+\lo) -- (8*\dx,\Yzero+\lo)
    -- (8*\dx,\Yzero+\hi) -- (9*\dx,\Yzero+\hi)
    -- (9*\dx,\Yzero+\lo) -- (10*\dx,\Yzero+\lo)
    -- (10*\dx,\Yzero+\hi) -- (11*\dx,\Yzero+\hi)
    -- (11*\dx,\Yzero+\lo) -- (12*\dx,\Yzero+\lo)
    -- (12*\dx,\Yzero+\hi) -- (13*\dx,\Yzero+\hi)
    -- (13*\dx,\Yzero+\lo) -- (14*\dx,\Yzero+\lo)
    -- (14*\dx,\Yzero+\hi) -- (15*\dx,\Yzero+\hi)
    -- (15*\dx,\Yzero+\lo) -- (16*\dx,\Yzero+\lo);
  \node[blue, font=\small] at (\xE+0.5, \Yzero+0.25) {Q0};

  % ====== Q1 (red): 0,0,1,1,0,0,1,1,... 每2拍翻转 ======
  \draw[red, very thick]
    (0,\Yone+\lo) -- (0,\Yone+\lo) -- (1*\dx,\Yone+\lo)
    -- (1*\dx,\Yone+\lo) -- (2*\dx,\Yone+\lo)
    -- (2*\dx,\Yone+\hi) -- (3*\dx,\Yone+\hi)
    -- (3*\dx,\Yone+\hi) -- (4*\dx,\Yone+\hi)
    -- (4*\dx,\Yone+\lo) -- (5*\dx,\Yone+\lo)
    -- (5*\dx,\Yone+\lo) -- (6*\dx,\Yone+\lo)
    -- (6*\dx,\Yone+\hi) -- (7*\dx,\Yone+\hi)
    -- (7*\dx,\Yone+\hi) -- (8*\dx,\Yone+\hi)
    -- (8*\dx,\Yone+\lo) -- (9*\dx,\Yone+\lo)
    -- (9*\dx,\Yone+\lo) -- (10*\dx,\Yone+\lo)
    -- (10*\dx,\Yone+\hi) -- (11*\dx,\Yone+\hi)
    -- (11*\dx,\Yone+\hi) -- (12*\dx,\Yone+\hi)
    -- (12*\dx,\Yone+\lo) -- (13*\dx,\Yone+\lo)
    -- (13*\dx,\Yone+\lo) -- (14*\dx,\Yone+\lo)
    -- (14*\dx,\Yone+\hi) -- (15*\dx,\Yone+\hi)
    -- (15*\dx,\Yone+\hi) -- (16*\dx,\Yone+\hi);
  \node[red, font=\small] at (\xE+0.5, \Yone+0.25) {Q1};

  % ====== Q2 (green): 0,0,0,0,1,1,1,1,... 每4拍翻转 ======
  \draw[green!50!black, very thick]
    (0,\Ytwo+\lo) -- (3*\dx,\Ytwo+\lo)
    -- (3*\dx,\Ytwo+\lo) -- (4*\dx,\Ytwo+\lo)
    -- (4*\dx,\Ytwo+\hi) -- (7*\dx,\Ytwo+\hi)
    -- (7*\dx,\Ytwo+\hi) -- (8*\dx,\Ytwo+\hi)
    -- (8*\dx,\Ytwo+\lo) -- (11*\dx,\Ytwo+\lo)
    -- (11*\dx,\Ytwo+\lo) -- (12*\dx,\Ytwo+\lo)
    -- (12*\dx,\Ytwo+\hi) -- (15*\dx,\Ytwo+\hi)
    -- (15*\dx,\Ytwo+\hi) -- (16*\dx,\Ytwo+\hi);
  \node[green!50!black, font=\small] at (\xE+0.5, \Ytwo+0.25) {Q2};

  % ====== Q3 (violet): 0,0,0,0,0,0,0,0,1,1,1,1,1,1,1,1 每8拍翻转 ======
  \draw[violet, very thick]
    (0,\Ythree+\lo) -- (7*\dx,\Ythree+\lo)
    -- (7*\dx,\Ythree+\lo) -- (8*\dx,\Ythree+\lo)
    -- (8*\dx,\Ythree+\hi) -- (15*\dx,\Ythree+\hi)
    -- (15*\dx,\Ythree+\hi) -- (16*\dx,\Ythree+\hi);
  \node[violet, font=\small] at (\xE+0.5, \Ythree+0.25) {Q3};

  % === 图例 ===
  \node[font=\footnotesize, align=left, anchor=north west]
    at (0, \Ythree+\hi+0.5) {
    \textcolor{blue}{Q0}: $f_{clk}/2$\quad
    \textcolor{red}{Q1}: $f_{clk}/4$\quad
    \textcolor{green!50!black}{Q2}: $f_{clk}/8$\quad
    \textcolor{violet}{Q3}: $f_{clk}/16$
  };
\end{tikzpicture}
\end{center}

\textbf{关键观察}：
\begin{itemize}
  \item Q0在\textbf{每个}时钟上升沿翻转——这是二分频的基础
  \item Q1每2个时钟周期翻转一次——Q0的翻转频率是Q1的2倍
  \item Q2每4个时钟周期翻转一次——Q1的翻转频率是Q2的2倍
  \item Q3每8个时钟周期翻转一次——Q2的翻转频率是Q3的2倍
  \item 规律：$Q_k$的频率 = $f_{clk} / 2^{k+1}$，这正是\textbf{分频器}的核心原理
\end{itemize}

\section{Testbench}

\begin{lstlisting}[caption={4位计数器Testbench}]
library ieee;
use ieee.std_logic_1164.all;

entity counter_4bit_tb is
end counter_4bit_tb;

architecture test of counter_4bit_tb is
  signal clk, rst_n, carry : std_logic;
  signal count : std_logic_vector(3 downto 0);
  constant clk_period : time := 20 ns;

  component counter_4bit is
    port (clk, rst_n : in std_logic;
          count : out std_logic_vector(3 downto 0);
          carry : out std_logic);
  end component;
begin
  uut: counter_4bit port map (clk, rst_n, count, carry);

  clk_process: process
  begin
    clk <= '0'; wait for clk_period/2;
    clk <= '1'; wait for clk_period/2;
  end process;

  stimulus: process
  begin
    rst_n <= '0'; wait for 30 ns;
    rst_n <= '1'; wait for 500 ns;  -- 观察完整16拍循环
    wait;
  end process;
end test;
\end{lstlisting}
